\section{Four Years Worked End to End}
\label{sec:years}

The rules above are stated at their loci; this section applies them. Four civil years are worked
through: 1962, the year the typical edition appeared and a year in which nothing is resumed and
nothing dropped; 2008, the extreme year in which only one Sunday after the Epiphany occurs and
twenty-eight Sundays after Pentecost do; and 2026 and 2027, two ordinary recent years with
different resumption cases. The rare shortfall year 2011 is noted at the end because
Section~\ref{sec:resumed} treats it as unresolved.

Every date below was computed from the Gregorian computus by the method of
Section~\ref{sec:movable}. For 1962, 2008, and 2011 the anchors were additionally compared with
the \latin{Tabella temporaria} printed in the Missal itself and agree in every column. For 2026
and 2027 no printed witness was consulted, and the dates are computed only. Nothing here is an
Ordo: a real celebration also depends on the diocesan, religious, and church calendar, on
patronal and dedication feasts, and on the actual decisions of the competent authority.

\subsection{1962: $E = 6$, $P = 24$, nothing resumed}

\begingroup\small
\renewcommand{\arraystretch}{1.08}
\begin{longtable}{@{}>{\raggedright\arraybackslash}p{0.36\linewidth} >{\raggedright\arraybackslash}p{0.26\linewidth} >{\raggedright\arraybackslash}p{0.30\linewidth}@{}}
\toprule
\textbf{Anchor} & \textbf{Date} & \textbf{Note} \\
\midrule
\endfirsthead
\toprule
\textbf{Anchor} & \textbf{Date} & \textbf{Note} \\
\midrule
\endhead
\bottomrule
\endfoot
Epiphany & Saturday 6 January & \\
Sundays after the Epiphany & 7, 14, 21, 28 January; 4, 11 February & All six occur; $E = 6$ \\
Most Holy Name of Jesus & Tuesday 2 January & No Sunday falls 2--5 January \\
Septuagesima & 18 February & \\
Ash Wednesday & 7 March & \\
Lenten Ember days & 14, 16, 17 March & \\
First Sunday of the Passion & 8 April & \\
Palm Sunday & 15 April & \\
Easter & 22 April & \\
Rogation days & 28--30 May & As of themselves (\rn{87}) \\
Ascension & Thursday 31 May & \\
Pentecost & 10 June & \\
Trinity Sunday & 17 June & First Sunday after Pentecost \\
Corpus Christi & Thursday 21 June & \\
Most Sacred Heart & Friday 29 June & \\
September Ember days & 19, 21, 22 September & Third Sunday of September is 16 September \\
Christ the King & 28 October & Last Sunday of October \\
Twenty-third Sunday after Pentecost & 18 November & \\
Last Sunday after Pentecost & 25 November & The twenty-fourth; $P = 24$ \\
First Sunday of Advent & 2 December & \\
Advent Ember days & 19, 21, 22 December & Third Sunday of Advent is 16 December \\
\end{longtable}
\endgroup


\textbf{Resumption.} $P = 24$ is neither \latin{plures quam XXIV} nor a shortfall. The
twenty-four printed formularies after Pentecost are used in order and no Epiphany Mass is
resumed. Since $E = 6$, none was left over either: $S = (6-6) - (24-24) = 0$. In 1962 the two
counts balance exactly, which happens in 64 of the 201 years from 1900 to 2100.

\textbf{Occurrences worth stating.}

\begin{itemize}[leftmargin=1.3em, itemsep=0.25em]
\item \textbf{19 March}, St Joseph (I class), falls on the Monday after the Second Sunday of Lent
— Ash Wednesday being 7 March, the Sundays of Lent are 11, 18, and 25 March.
The day is a Lenten feria of the III class (\rn{25}\,a); the feast holds the higher place at line
11, so the feast is celebrated. Because the feria is of Lent, its commemoration is privileged
(\rn{109}\,e), and \rn{111}\,a admits one privileged commemoration on a I class day; the feria is
therefore commemorated.
\item \textbf{25 March}, the Annunciation (I class), falls on the Third Sunday of Lent. Sundays
of Lent stand at line 6, above I class feasts of the universal Church at line 11, so the Sunday
prevails. The feast, being of the I class, has the right of translation under \rn{95} and goes,
by \rn{96}, to the next following day that is not of the I or II class — Monday 26 March, a
Lenten feria of the III class. The special seat of \rn{96}\,a, the Monday after Low Sunday, is
not engaged, because it applies only when the Annunciation must be transferred \emph{after
Easter}.
\item \textbf{21 December}, St Thomas the Apostle (II class of the universal Church, line 16),
falls on the Ember Friday of Advent (II class feria, line 18). The feast prevails; the Ember
feria, being of the II class, must be commemorated when impeded (\rn{24}).
\item \textbf{22 December} is the Ember Saturday of Advent. By Missal rubrics \rn{300} the Mass
in which sacred Orders are conferred that day must be of the Saturday, even if a I or II class
feast occurs.
\end{itemize}

\subsection{2008: $E = 1$, $P = 28$, the maximum resumption}

2008 is the more instructive of the two extremes. The Epiphany fell on a Sunday, Septuagesima
came on 20 January, and only one Sunday after the Epiphany occurred; Easter on 23 March is close
to the earliest possible, and the year carried twenty-eight Sundays after Pentecost.

\begingroup\small
\renewcommand{\arraystretch}{1.08}
\begin{longtable}{@{}>{\raggedright\arraybackslash}p{0.36\linewidth} >{\raggedright\arraybackslash}p{0.26\linewidth} >{\raggedright\arraybackslash}p{0.30\linewidth}@{}}
\toprule
\textbf{Anchor} & \textbf{Date} & \textbf{Note} \\
\midrule
\endfirsthead
\toprule
\textbf{Anchor} & \textbf{Date} & \textbf{Note} \\
\midrule
\endhead
\bottomrule
\endfoot
Most Holy Name of Jesus & Wednesday 2 January & No Sunday falls 2--5 January \\
Epiphany & Sunday 6 January & \\
Sundays after the Epiphany & 13 January only & $E = 1$ \\
Septuagesima & 20 January & \\
Ash Wednesday & 6 February & \\
Lenten Ember days & 13, 15, 16 February & \\
Palm Sunday & 16 March & \\
Easter & 23 March & \\
Rogation days & 28--30 April & \\
Ascension & Thursday 1 May & \\
Pentecost & 11 May & \\
Trinity Sunday & 18 May & \\
Corpus Christi & Thursday 22 May & \\
September Ember days & 24, 26, 27 September & Third Sunday of September is 21 September \\
Christ the King & 26 October & \\
Twenty-third Sunday after Pentecost & 19 October & \\
Last Sunday after Pentecost & 23 November & The twenty-eighth in the count; $P = 28$ \\
First Sunday of Advent & 30 November & \\
Advent Ember days & 17, 19, 20 December & \\
\end{longtable}
\endgroup


\textbf{13 January: a collision the Missal answers by name.} The First Sunday after the Epiphany
carries the Holy Family (II class) by \rn{17}\,b; 13 January carries the Commemoration of the
Baptism of the Lord (II class). The Missal prints the resolution twice, once at each formulary:
\latin{Si festum S. Familiae occurrerit die 13 ianuarii, Missa dicitur de festo S. Familiae, sine
commemoratione Baptismatis D. N. I. C., et sine commemoratione dominicae} — if the feast of the
Holy Family occurs on 13 January, the Mass is said of the feast of the Holy Family, without a
commemoration of the Baptism of Our Lord Jesus Christ and without a commemoration of the Sunday.
Both commemorations that the general rules would otherwise generate are expressly suppressed.

\textbf{Resumption.} With $E = 1$, five Epiphany formularies were unused: the Second through the
Sixth. With $P = 28$, case \textit{d} of \rn{18} fills four places. The assignment is therefore:

\begin{center}
\small
\begin{tabular}{@{}r l >{\raggedright\arraybackslash}p{0.46\linewidth}@{}}
\toprule
\textbf{Place} & \textbf{Date} & \textbf{Formulary} \\
\midrule
23 & 19 October & Twenty-third Sunday after Pentecost \\
24 & 26 October & Third Sunday left over after the Epiphany \\
25 & 2 November & Fourth Sunday left over after the Epiphany \\
26 & 9 November & Fifth Sunday left over after the Epiphany \\
27 & 16 November & Sixth Sunday left over after the Epiphany \\
28 & 23 November & \latin{Dominica XXIV et ultima post Pentecosten} \\
\bottomrule
\end{tabular}
\end{center}

The Second Sunday after the Epiphany occurred nowhere in 2008. It was impeded in January by
Septuagesima and found no resumption slot in November, because \rn{18} fills the slots from the
Sixth downward and case \textit{d} reaches only the Third. This is the omission that \rn{18}'s
closing clause contemplates, and 2008 is one of only two years between 1900 and 2100 in which it
falls on the Second rather than on a higher ordinal.

\textbf{Two of the resumed places were themselves impeded.}

\begin{itemize}[leftmargin=1.3em, itemsep=0.25em]
\item \textbf{26 October} was the last Sunday of October and therefore Christ the King (I class),
which by \rn{17}\,d is perpetually assigned to that Sunday and \latin{locum tenet dominicae
occurrentis cum omnibus iuribus et privilegiis}, no commemoration of the Sunday being made. The
Third Sunday left over after the Epiphany, assigned to the twenty-fourth place, therefore did
not occur either. Its assignment by \rn{18} is an identification of the Sunday, not a guarantee
that its Mass is said.
\item \textbf{2 November} was a Sunday, so the Commemoration of All the Faithful Departed, which
at line 8 of the table expressly \latin{locum cedit dominicae occurrenti}, yielded and was
transferred by \rn{96}\,b, \latin{tamquam in sedem propriam}, to Monday 3 November. The Fourth
Sunday left over after the Epiphany was celebrated on 2 November.
\end{itemize}

\textbf{9 November, a case this reference does not settle.} The twenty-sixth place fell on the
Dedication of the Archbasilica of the Most Holy Saviour, a feast of the II class in the universal
calendar. If that feast is reckoned among the \latin{festa Domini}, it stands at line 14, above
Sundays of the II class at line 15, and by \rn{16}\,a it would take the Sunday's place with no
commemoration of the Sunday. If it is not, it stands at line 16, below the Sunday, and by
\rn{111}\,b the Sunday admits its commemoration. The 1962 calendar page does not mark feasts of
the Lord as such, and this reference did not establish the classification from a controlling
text; the competent Ordo states it. The point is recorded here rather than resolved, because a
reference that guesses at it will mislead a reader in every year in which the collision recurs.

\subsection{2026: $E = 3$, $P = 26$, two resumed and one dropped}

\begingroup\small
\renewcommand{\arraystretch}{1.08}
\begin{longtable}{@{}>{\raggedright\arraybackslash}p{0.36\linewidth} >{\raggedright\arraybackslash}p{0.26\linewidth} >{\raggedright\arraybackslash}p{0.30\linewidth}@{}}
\toprule
\textbf{Anchor} & \textbf{Date} & \textbf{Note} \\
\midrule
\endfirsthead
\toprule
\textbf{Anchor} & \textbf{Date} & \textbf{Note} \\
\midrule
\endhead
\bottomrule
\endfoot
Most Holy Name of Jesus & Sunday 4 January & Falls within 2--5 January \\
Epiphany & Tuesday 6 January & \\
Sundays after the Epiphany & 11, 18, 25 January & $E = 3$ \\
Septuagesima & 1 February & \\
Ash Wednesday & 18 February & \\
Lenten Ember days & 25, 27, 28 February & \\
First Sunday of the Passion & 22 March & \\
Palm Sunday & 29 March & \\
Easter & 5 April & \\
Rogation days & 11--13 May & \\
Ascension & Thursday 14 May & \\
Pentecost & 24 May & \\
Trinity Sunday & 31 May & \\
Corpus Christi & Thursday 4 June & \\
Most Sacred Heart & Friday 12 June & \\
September Ember days & 23, 25, 26 September & Third Sunday of September is 20 September \\
Christ the King & 25 October & \\
First Sunday of Advent & 29 November & \\
Advent Ember days & 16, 18, 19 December & \\
\end{longtable}
\endgroup


\textbf{Resumption.} Three Epiphany formularies were unused — the Fourth, Fifth, and Sixth — and
$P = 26$ gives two places. Case \textit{b} of \rn{18} assigns the Fifth to the twenty-fourth
place and the Sixth to the twenty-fifth; the Fourth is dropped. $S = (6-3) - (26-24) = 1$.

\begin{center}
\small
\begin{tabular}{@{}r l >{\raggedright\arraybackslash}p{0.52\linewidth}@{}}
\toprule
\textbf{Place} & \textbf{Date} & \textbf{Formulary} \\
\midrule
23 & 1 November & Twenty-third Sunday after Pentecost — but see below \\
24 & 8 November & Fifth Sunday left over after the Epiphany \\
25 & 15 November & Sixth Sunday left over after the Epiphany \\
26 & 22 November & \latin{Dominica XXIV et ultima post Pentecosten} \\
\bottomrule
\end{tabular}
\end{center}

\textbf{1 November.} All Saints, a I class feast of the universal Church, falls on the Sunday
occupying the twenty-third place. It stands at line 11 and the Sunday at line 15, so the feast is
celebrated. All Saints is not a feast of the Lord, so the substitution rule of \rn{16}\,a does
not apply and the Sunday is not simply absorbed; it is commemorated, its commemoration being
privileged under \rn{109}\,a and permitted by \rn{111}\,a as the one privileged commemoration
allowed on a I class day. The Mass of the Twenty-third Sunday after Pentecost is therefore not
said in 2026, though the Sunday is commemorated.

\textbf{Other occurrences.} 19 March falls on the Thursday after the Fourth Sunday of Lent, a
Lenten feria of the III class; St Joseph prevails and the feria is commemorated. 25 March falls
on the Wednesday of Passion week, likewise a III class feria under \rn{25}\,a; the Annunciation
prevails and the feria is commemorated. Neither feast is transferred, because neither meets a
day standing higher in the table.

\subsection{2027: $E = 2$, $P = 27$, three resumed and one dropped}

\begingroup\small
\renewcommand{\arraystretch}{1.08}
\begin{longtable}{@{}>{\raggedright\arraybackslash}p{0.36\linewidth} >{\raggedright\arraybackslash}p{0.26\linewidth} >{\raggedright\arraybackslash}p{0.30\linewidth}@{}}
\toprule
\textbf{Anchor} & \textbf{Date} & \textbf{Note} \\
\midrule
\endfirsthead
\toprule
\textbf{Anchor} & \textbf{Date} & \textbf{Note} \\
\midrule
\endhead
\bottomrule
\endfoot
Most Holy Name of Jesus & Sunday 3 January & \\
Epiphany & Wednesday 6 January & \\
Sundays after the Epiphany & 10, 17 January & $E = 2$ \\
Septuagesima & 24 January & \\
Ash Wednesday & 10 February & \\
Lenten Ember days & 17, 19, 20 February & \\
First Sunday of the Passion & 14 March & \\
Palm Sunday & 21 March & \\
Easter & 28 March & \\
Rogation days & 3--5 May & \\
Ascension & Thursday 6 May & \\
Pentecost & 16 May & \\
Trinity Sunday & 23 May & \\
Corpus Christi & Thursday 27 May & \\
Most Sacred Heart & Friday 4 June & \\
September Ember days & 22, 24, 25 September & Third Sunday of September is 19 September \\
Christ the King & 31 October & \\
First Sunday of Advent & 28 November & \\
Advent Ember days & 15, 17, 18 December & \\
\end{longtable}
\endgroup


\textbf{Resumption.} Four Epiphany formularies were unused — the Third through the Sixth — and
$P = 27$ gives three places. Case \textit{c} assigns the Fourth, Fifth, and Sixth to places 24,
25, and 26; the Third is dropped. $S = (6-2) - (27-24) = 1$.

\begin{center}
\small
\begin{tabular}{@{}r l >{\raggedright\arraybackslash}p{0.52\linewidth}@{}}
\toprule
\textbf{Place} & \textbf{Date} & \textbf{Formulary} \\
\midrule
23 & 24 October & Twenty-third Sunday after Pentecost \\
24 & 31 October & Fourth Sunday left over after the Epiphany — impeded by Christ the King \\
25 & 7 November & Fifth Sunday left over after the Epiphany \\
26 & 14 November & Sixth Sunday left over after the Epiphany \\
27 & 21 November & \latin{Dominica XXIV et ultima post Pentecosten} \\
\bottomrule
\end{tabular}
\end{center}

As in 2008, the last Sunday of October falls in a resumed place, and Christ the King takes it
under \rn{17}\,d without commemoration of the Sunday. The Fourth Sunday left over after the
Epiphany therefore does not occur in 2027 either, having been assigned a place that a
perpetually-assigned I class feast occupies.

\textbf{25 April: the greater Litanies on a Sunday.} In 2027 the greater Litanies fall on the
Fourth Sunday after Easter. The transfer clause of \rn{80} is not engaged, because it operates
only when Easter Sunday or Easter Monday falls on 25 April. The layers are:

\begin{enumerate}[leftmargin=1.6em, itemsep=0.2em]
\item The day is a Sunday of the II class (line 15). St Mark the Evangelist, a II class feast of
the universal Church that is not of the Lord, stands at line 16 and yields.
\item \rn{111}\,b admits on a II class Sunday one commemoration only, that of a II class feast,
\latin{quae tamen omittitur si commemoratio privilegiata facienda sit} — which is omitted if a
privileged commemoration must be made. The commemoration of the greater Litanies is privileged in
the Mass by \rn{109}\,f.
\item In churches where the procession or the supplications appointed by the Ordinary under
\rn{83} are held, the Rogation Mass is said as a votive Mass of the II class (Missal rubrics
\rnn{341}, 346), regularly after the procession (\rn{347}); where that votive Mass is impeded,
its oration is added to the day's Mass under a single conclusion (\rn{343}\,c), the day not being
one of those listed at numbers 1, 2, 3, and 8 of the table.
\item \rn{81} adds that the commemoration of the greater Litanies is not to be counted a
commemoration ``of the season'', so that the bar of \rn{112}\,c against a second seasonal
commemoration is not triggered.
\end{enumerate}

Which of these dispositions actually obtains in a given church depends on whether the procession
or supplications are held there, which \rn{82} leaves to the local Ordinary. That is precisely
the kind of question a computed calendar cannot answer and an Ordo can.

\subsection{2011, and why it is left open}

2011 is one of the four years between 1900 and 2100 with twenty-three Sundays after Pentecost.
Easter fell on 24 April, Pentecost on 12 June, and the First Sunday of Advent on 27 November, so
that the last Sunday after Pentecost was 20 November and it was the twenty-third in the count.
All six Sundays after the Epiphany occurred, so nothing was left over on that side. By the
arithmetic of Section~\ref{sec:movable}, $S = (6-6) - (23-24) = 1$: one printed formulary too
many, and both candidates for omission — the Twenty-third and the Twenty-fourth — stand in the
after-Pentecost series. Section~\ref{sec:resumed} sets out the two readings and states what
would settle the question. This reference does not choose between them, and a reader preparing
for 2038 or 2095 should consult a competent Ordo rather than this document on that single point.

One other feature of 2011 deserves mention because it exercises Missal rubrics \rn{300}: the
Ember Saturday of Lent fell on 19 March, the feast of St Joseph. The I class feast prevails over
the II class Ember feria in occurrence, and the feria must be commemorated under \rn{24}; but the
Mass in which sacred Orders are conferred that day is still of the Saturday, \latin{etiam festo I
vel II classis occurrente}.
